\subsection{Hipóteses de trabalho e critérios de verificação}

As hipóteses assumem a forma de respostas provisórias verificáveis, que poderão ser corroboradas, qualificadas ou rejeitadas pelas evidências empíricas. Elas não configuram um teste causal clássico entre uma variável independente e uma variável dependente única; orientam a bateria de perguntas, a seleção dos critérios avaliativos e a interpretação dos resultados no desenho orientado a artefato.

\begin{enumerate}[label=\textbf{H\arabic*.}]
  \item \textbf{A pertinência das evidências recuperadas é condição necessária, mas não suficiente, para a verificabilidade das respostas.} A hipótese será examinada tanto nos casos sem evidência pertinente para sustentar as afirmações quanto naqueles em que existe evidência pertinente, mas a geração não a utiliza adequadamente ou extrapola o contexto recuperado. Será corroborada se os dois padrões ocorrerem em casos documentados; qualificada se dependerem do tipo de pergunta ou de condições específicas; e rejeitada se a verificabilidade for integralmente explicada pela qualidade da recuperação.

  \item \textbf{O desempenho tende a ser inferior em perguntas que exigem integração de múltiplas evidências, raciocínio em etapas ou reconhecimento de informação ausente, em comparação com perguntas factuais simples, delimitadas por autor e apoiadas em evidência direta.} A hipótese será examinada por comparação de desempenho e análise de erros entre categorias. Será corroborada se as categorias mais complexas apresentarem desempenho consistentemente inferior; qualificada se a diferença se restringir a dimensões ou condições específicas; e rejeitada se o desempenho for homogêneo ou inverso ao esperado \parencite{saadFalconEtAl2024ares,esEtAl2023ragas,zhuEtAl2025rageval,pipitoneAlami2024legalbenchRag}.

  \item \textbf{A avaliação baseada no paradigma \textit{LLM-as-a-judge} apresentará concordância parcial com a avaliação humana, mas será insuficiente como critério único de qualidade.} A hipótese será examinada pela comparação com julgamentos humanos independentes. Será corroborada se houver concordância parcial acompanhada de divergências materialmente relevantes; qualificada se a concordância variar por tipo de pergunta, dimensão ou rubrica; e rejeitada se os julgamentos automatizados forem plenamente consistentes com os humanos ou incapazes de discriminar qualidade \parencite{liuEtAl2023geval,liuEtAl2025judgeRag}.
\end{enumerate}
